\subsection{互素元}

定义 5. --- 在序群中，两个元素 x 与 y 称为互素的，如果 \(\inf(x, y) = 0\) 。

在某些情形下，把满足 \(\inf(|x|, |y|) = 0\) 的两个元素定义为互素是自然的（参见 INT，II，§ 1），或者在整除理论中引入相应的术语。我们不这样做。

两个互素元素必定是正的。x 的正部和负部 \(x^{+}\) 与 \(x^{-}\) 是互素的（VI，第12页，命题9 a））。元素 \(x_{i}\) （属于某个族 \((x_{t})_{t\in I}\) ）称为整体互素的，如果 \(\inf_{t\in I}x_t = 0\) ；若这些 \(x_{t}\) 都 \(\geqslant 0\) ，则

只需存在 I 的一个有限子集 J，使得相应的元素整体互素。族 \((x_{\iota})\) 的元素称为两两互素的，如果 \(\inf(x_{\iota}, x_{\kappa}) = 0\) 对每对不同的指标 \((\iota, \kappa)\) 成立。

诸 \(x_{i}\) 可以整体互素而不两两互素。

若 \(x\) 与 \(y\) 是互素的，我们也说 \(x\) 与 \(y\) 互素，或说 \(y\) 与 \(x\) 互素。

（DIV）K 的两个元素 x 与 y 称为互素的，如果主理想 \((x)\) 与 \((y)\) 在 \(P^{*}\) 中非零且互素；这等于说 1 是 x 与 y 的一个最大公因子，并蕴含 x 与 y 属于 A。例如，既约分数的分子与分母是互素的。两两互素元素和整体互素元素的概念可类似地定义。

（DIV）当 \(x\) 与 \(y\) 互素时，常称它们“彼此互素”；宜避免这一术语，它可能与素数的概念（I，第50页，定义16）混淆。

命题 10. --- 设 \(x, y\) 与 \(z\) 是序群的三个元素；要使 \(x - z\) 与 \(y - z\) 互素，必要且充分的是 \(z = \inf(x, y)\) 。

事实上，关系 \(z = \inf (x, y)\) 与 \(0 = \inf (x - z, y - z)\) 是等价的。

命题10（DIV）。 --- 设 \(a, b\) 与 \(c\) 是 \(\mathbf{K}\) 的三个元素，且 \(c \neq 0\) ；要使商 \(ac^{-1}\) 与 \(bc^{-1}\) 互素，必要且充分的是 \(c\) 是 \(a\) 与 \(b\) 的一个最大公因子。

命题11. --- 若 \((x_{i})\) , \((y_{j})\) 是格序群中正元素的两个有限族，则

\[
\inf \left(\sum_ {i} x _ {i}, \sum_ {j} y _ {j}\right) \leqslant \sum_ {i, j} \inf (x _ {i}, y _ {j}).
\]

对族 \((x_{i})\) 与 \((y_{j})\) 的元素个数进行归纳论证，只需证明：如果 \(x, y\) 与 \(z\) 是正元素，则

\[
\inf (x, y + z) \leqslant \inf (x, y) + \inf (x, z).
\]

事实上，令 \(t = \inf(x, y + z)\) ；由VI第12页的推论，我们可以写 \(t = t_1 + t_2\) ，其中 \(0 \leqslant t_1 \leqslant y\) 且 \(0 \leqslant t_2 \leqslant z\) ；由于 \(t_1\) 与 \(t_2\) 为正，我们还有 \(t_1 \leqslant x\) 与 \(t_2 \leqslant x\) ，从而 \(t_1 \leqslant \inf(x, y)\) 与 \(t_2 \leqslant \inf(x, z)\) 。

推论1. --- 若x与y是格序群的两个互素元，且z是一个正元素，则 \(\inf(x, z) = \inf(x, y + z)\) 。

事实上， \(\inf (x,y + z)\leqslant \inf (x,z)\) 由命题11得出，且 \(\inf (x,z)\leqslant \inf (x,y + z)\) 由 \(y\geqslant 0\) 得出，由此即得推论。

推论2. --- 在格序群中，若x与y互素，且 \(z \geqslant 0\) 与 \(x \leqslant y + z\) ，则 \(x \leqslant z\) 。

推论3. --- 在一个格序群中，如果x与y互素且与z互素，则x也与 \(y + z\) 互素。

推论4. --- 如果 \((x_i)_{1 \leq i \leq n}\) 与 \((y_j)_{1 \leq j \leq m}\) 是格序群G中元素的两个有限族，使得每个 \(x_i\) 与每个 \(y_j\) 互素，则 \(x_1 + \cdots + x_n\) 与 \(y_1 + \cdots + y_m\) 互素。

这由推论3通过对m和n进行归纳得出。

推论5. --- 对任意整数 \(n \geqslant 0\) ，我们有 \((nx)^+ = nx^+\) 与 \((nx)^- = nx^-\) ；对所有 \(n \in \mathbf{Z}\) ，我们有 \(|nx| = |n| \cdot |x|\) 。

事实上， \(nx = nx^{+} - nx^{-}\) ；由于 \(x^{+}\) 与 \(x^{-}\) 互素，因此 \(nx^{+}\) 与 \(nx^{-}\) 在 \(n \geqslant 0\) 时也互素（推论4）；第一个断言由VI第12页命题9 b)得出。第二个断言在情形 \(n \geqslant 0\) 下由第一个断言根据命题9 d)得出；情形 \(n < 0\) 由此根据关系 \(|-x| = |x|\) 得出。

命题11（DIV）。--- 假设集合 \(\mathcal{P}^*\) 是一个格，并设 \((a_i)\) , \((b_j)\) 为A中元素的两个有限族。则 \(\prod_{i} a_i\) 与

\(\prod_{j} b_{j}\) 的每一个最大公因子都整除乘积 \(\prod_{i,j} \gcd(a_{i}, b_{j})\) 。

推论1（DIV）。--- 若 a、b、c 是A的三个元素，且a与b互素，则a与c的每个最大公因子也是a与bc的一个最大公因子。

推论2（DIV）（欧几里得引理）。--- 设 \(a, b, c\) 是A的三个元素。如果 \(a\) 与 \(b\) 互素且整除 \(bc\) ，则它整除 \(c\) 。

推论3（DIV）。--- 若x与y互素且与z互素，则x与yz互素。

推论4（DIV）。--- 如果 \((x_{i})\) 与 \((y_{j})\) 是A中元素的两个有限族，使得每个 \(x_{i}\) 与每个 \(y_{j}\) 互素，则这些 \(x_{i}\) 的乘积与这些 \(y_{j}\) 的乘积互素。

推论5（DIV）。--- 如果 \(d\) 是 \(x\) 与 \(y\) 的一个最大公因子，则 \(d^n\) 是 \(x^n\) 与 \(y^n\) 的一个最大公因子（对每个正整数 \(n\) ）。

事实上， \(xd^{-1}\) 与 \(yd^{-1}\) 互素（命题10（DIV）），因此 \(x^{n}d^{-n}\) 与 \(y^{n}d^{-n}\) 也互素（推论4）。

命题12. --- 设 \(x_{i}\) ( \(1 \leqslant i \leqslant n\) ) 是格序群中 \(n\) 个两两互素的元素。则

\[
\sup \left(x _ {1}, \dots , x _ {n}\right) = x _ {1} + \dots + x _ {n}.
\]

这由公式 \(u + v = \sup(u, v) + \inf(u, v)\) （VI，第10页，命题7）通过对n归纳得出，其中用到 \(x_{i}\) 与 \(x_{1} + \cdots + x_{i-1}\) 互素这一事实， \(2 \leqslant i \leqslant n\) （命题11的推论4）。

注. --- VI第10页命题7还表明， \(x\) 与 \(y\) 互素的一个充分必要条件是 \(x + y = \sup(x, y)\) 。

命题12（DIV）。--- 设 \(a_i\) 是 \(n\) 个属于 \(A\) 的两两互素的元素；则乘积 \(a_1 \ldots a_n\) 是 \(a_1, \ldots, a_n\) 的一个最小公倍元。

命题13. --- 在格序群G中，设 \((x_{\alpha})\) 是一个具有下确界（相应地，上确界）的族，并设 z 为 G 的任意元素；则族 \((\sup(z, x_{\alpha}))\) （相应地 \((\inf(z, x_{\alpha})))\) 具有下确界（相应地，上确界）且

\[
\left\{ \begin{array}{l} \inf _ {\alpha} (\sup (z, x _ {\alpha})) = \sup \left(z, \inf _ {\alpha} x _ {\alpha}\right) \\ \sup _ {\alpha} (\inf (z, x _ {\alpha})) = \inf \left(z, \sup _ {\alpha} x _ {\alpha}\right) \end{array} \right.\tag{4}
\]

（相应地）。

假设族 \((x_{\alpha})\) 具有下确界 \(y\) 。我们将证明 \(\sup (z,y)\) 是族 \((\sup (z,x_{\alpha}))\) 的一个下确界。

事实上， \(\sup (z, x_{\alpha}) = z + (x_{\alpha} - z)^{+}\) ，通过平移我们可以归结为 \(z = 0\) 的情形，即我们必须证明 \((x_{\alpha}^{+})\) 的下确界等于 \(y^{+}\) 。由于 \(y \leqslant x_{\alpha}\) ，我们有 \(y^{+} \leqslant x_{\alpha}^{+}\) 对所有 \(\alpha\) 成立（VI，第12页，命题9 c））。反之，若 \(a \leqslant x_{\alpha}^{+}\) 对所有 \(\alpha\) ，则由此可得 \(a \leqslant x_{\alpha} + x_{\alpha}^{-}\) （命题9 a））；现在 \(y^{-} \geqslant x_{\alpha}^{-}\) ，因为 \(y \leqslant x_{\alpha}\) ；因此 \(a \leqslant x_{\alpha} + y^{-}\) 对所有 \(\alpha\) ，即 \(a \leqslant y + y^{-} = y^{+}\) 。

另一公式通过改为相反的序关系即可得出。

推论. --- 若格序群G的元素z与某个下确界为y的族的每个元素 \(x_{\alpha}\) 互素，则z与y互素。

这是公式 (4) 中第二个公式的直接推论。

注. --- 将命题13的公式应用于由两个元素 \((x, y)\) 组成的族，我们得到以下公式，它们表达了格序群中的两个合成律 sup、inf 各自对另一个的分配性：

\[
\begin{array}{l} \sup (z, \inf (x, y)) = \inf (\sup (z, x), \sup (z, y)) \\ \inf (z, \sup (x, y)) = \sup (\inf (z, x), \inf (z, y)). \end{array}
\]

这一分配性质是格序群所特有的，并不推广到任意格，甚至不推广到格序幺半群（参见VI，第33页，习题24）。
% semantic-fragments: legacy-input-seam
\relax{}
